%% notes-tex/database/database-expansion-hosts/sections/regdna.tex
%% [[experiments.database.expansion-hosts]]
%% https://github.com/Mjvolk3/torchcell/tree/main/notes-tex/database/database-expansion-hosts/sections/regdna.tex
%% SOURCE: hand-authored prose. Counts come from the triage report
%% /scratch/projects/torchcell-scratch/tmp/deep-research/reports/Pk and next dataset wave triage.md
%% and from .../research_notes/Pk and next dataset wave triage/scerevisiae_gene_level_readouts.md
%% and .../yeast_expression_fm_corpus.md; the served Vaishnav figures are read from
%% torchcell/datasets/scerevisiae/vaishnav2022.py.

\section{Regulatory DNA\secstatus{todo}}\label{sec:regdna}

The \term{sequence-to-readout} tier is the set of datasets whose perturbation is a specified
DNA sequence, meaning a promoter, a 5' UTR, a core promoter, a transcription start site, a 3'
end, a terminator, or an integration coordinate, and whose readout is a reporter level, a
translation output, or the fitness of the strain carrying it. The served graph holds one such
dataset at scale, \texttt{vaishnav2022.py} with 71,158 records, against 70 yeast datasets whose
perturbation space is a gene list. The fitness and gene-interaction maps alone supply 67
percent of all 65,889,109 yeast records, so adding to that axis adds to the two-thirds the
corpus already has, while this tier's perturbation space is continuous rather than enumerable,
which is the perturbation space a model of regulatory grammar is fitted on, and no deletion
collection supplies it.

\notesec{The de Boer and Vaishnav lineage} de Boer et al. 2020 measured the expression output
of 102,371,025 promoter sequences across five experiments. The design is an 80 bp fully random
insert at $-170$ to $-90$ relative to the transcription start site inside a fixed scaffold,
driving YFP against a constitutive RFP on an episomal vector, with cells sorted into 18 uniform
bins by the log ratio of the two reporters and the promoters in each bin sequenced, so the
expression estimate is a read-weighted bin average. The media are complex and there are three
of them, glucose, galactose, and glycerol, which makes the medium a genuine environment
variable rather than a single condition. Vaishnav et al. 2022 ran the same platform with two
training libraries, 30,722,376 random sequences in complex medium and 20,616,659 in defined
medium, and the graph already serves the Vaishnav child of those two. The DREAM 2022 release is
the third member at 6,739,258 training sequences plus 71,103 held out, and every published
yeast expression model is scored on that held-out set.

The per-record cost of the largest row is characterized only in aggregate and belongs in the
record as such. Because each library's complexity exceeded the number of sorted cells, 78
percent of de Boer's promoter sequences appear in only one bin, often as one read from one
cell, which the authors assess as roughly 24 percent error in the expression estimate on
held-out data. That is a typed gap or a read-count field rather than a reason to decline the
row, and \texttt{vaishnav2022.py} already carries the same information as \texttt{read\_count}.
Two provenance items are open: the paper names GEO GSE104903 and GSE104878 and never states
which accession holds which library, and it never itemizes the five experiments that sum to
102,371,025, so one GEO page check settles both.

\notesec{The superset admission the existing loader already makes possible} The served
Vaishnav store holds 71,158 records of \texttt{VariablePromoterReporterPerturbation} crossed
with \texttt{PooledSortSeqExpressionPhenotype}, of which 71,054 carry an estimate and 104
\texttt{seqID} groups carry a typed \texttt{ProvenanceGap} for verbatim \texttt{NA} cells. The
two deferred training libraries are the complex-medium and defined-medium training files on
Zenodo, 3.59 GB and 2.47 GB, plus a 0.67 GB table on GEO,
and the loader's own docstring already calls them ``a size decision for the owner,
not a schema or a provenance one.'' Admitting them is the superset case the incremental path
was built for: the served records are a proven subset, so the admission check can prove the
development store still produces every served experiment id plus more, and the increment
matches served nodes by content id and adds only the new records. Admitting de Boer needs no
new class either, since the same two classes serve its child. The DREAM release is the one
member of the three that needs more than a download, because it is a curated slice of the same
platform and therefore needs an \texttt{AggregationRecord} at G1 rather than entering as a
primary measurement.

\notesec{The 3' end and the position-effect rows} Shalem et al. 2015 measured 13,000 designed
3' end sequences, 102 bp variable inside a 150 bp oligo, by FACS-seq into 16 bins of a
two-reporter ratio, and the library sequences and their measured expression values sit together
in one supplementary XLSX with no accession. It extends sequence-level coverage from the
promoter to the terminator and sits among the eight benchmark tasks the species-aware DNA
language model is scored on, so ingesting it serves a dataset a published model is measured
against. Its only uncertainty figure is a library-wide technical relative standard deviation of
13.2 percent rather than a per-construct error. The position-effect rows are the
one-construct-many-loci shape, where the sequence variable is purely the integration
coordinate. There are four of them: a 1,044-locus panel with a greater-than-13-fold expression
range and location dominating promoter, reporter, and carbon source, a panel of about 250 loci
each with a full transcriptome rather than one reporter value, a 482-locus panel from an
earlier study, and a genome-wide scan reporting 428 high-expressed and 444 low-expressed
sites. Not one of the four has a located accession, and the largest states that its table is
available from the corresponding author on request, so the whole sub-tier is supplement-bound
and two of its rows are asks.

\notesec{What the rest of the tier holds} Four further rows carry $10^{4}$ measurements or
more each and each needs exactly one new leaf. A Pol II MASTER library measures initiation
efficiency at 1,000,000 transcription start sites, deposited at GEO GSE185290 and BioProject
PRJNA766624, and its own preprint reported about 80,000 variants against the published
1,000,000, so the release has to be pinned rather than the paper cited. A random 5' UTR library
detected 489,348 of more than 500,000 designed variants upstream of \gene{HIS3}, read out by
competitive growth rather than by sorting, which by the rule that a different technology is its
own leaf makes it a separate phenotype class from a sort-seq bin average. A promoter-swap
library places 130 synthetic promoters spanning a 500-fold range upstream of named native
genes, giving about 10,000 strains with a per-strain pooled fitness, where the released gene
count is reported as 81, 96, and about 100 in three places and has to be reconciled before a
loader is written. And a natural-variant MPRA measures 5,832 variants in the promoters of 2,503
genes with the tested 144 bp windows released, which is the only row found that pins a native
variant to a measured cis effect on a named gene. The terminator panels are the weakest tier
for machine readability: 100 terminators span 0.0613 to 1.8002 relative to the \gene{PGK1t}
standard and a second study characterized over 3,000 genomic terminators, and no repository
deposit was confirmed for any terminator panel, so all of them are supplement-only.

\notesec{The expression foundation models are not the candidates, and the libraries under them are}
Two 2025 yeast models assemble a corpus, and reading their corpora returns mostly data the
graph already serves. The fungal DNA language model's supervised corpus is 5,215 tracks, of
which 1,014 are the 1,000-isolate RNA-seq the graph serves as Caudal 2024, 1,128 ChIP-exo and
20 ChIP-MNase tracks come from one 2021 deposit at GEO GSE147927, and 3,053 are new induction
RNA-seq tracks generated at Calico. The single-cell model trains on expression matrices with no
DNA sequence anywhere, and six of its named sources already have loaders. So the models supply
a triage shortlist rather than a dataset, and their three non-served items are each retrievable
on their own terms. The 3,053 new tracks are the exception that is worth an ask rather than a
parse: they are distributed only as BigWig and TFRecord on a Google Cloud Storage bucket with
no GEO, SRA, Zenodo, or Hugging Face accession anywhere in the paper, so a sha256-pinned mirror
is the only reproducible route and there is no archive fallback if the bucket moves. Two
further rows are dropped from this tier rather than ranked in it, and both are recorded in
Table~\ref{tab:excluded}. The synthetic-chromosome phenotype claims appear as figures, and the
matching deposit lists 231 samples across genome sequencing, proteome, and transcriptome with
no phenotype data type at all. The aggregate deep-mutational-scanning route needs a live query
before it can be counted, because the curated aggregate holds 217 assays over 186 proteins and
its yeast-only subset has to be read off the live service after a second aggregate's search
endpoint returned 404.
